% !PRINTED 648
\weberpage{648}
%<<CONT
tion in lineare Faktoren, so muß nach 2. einer dieser Linearfaktoren durch $\pi$ teilbar sein. Es können nicht zwei verschiedene Linearfunktionen $z - c$ und $z - c'$ durch dasselbe $\pi$ teilbar sein, weil sonst die von Null verschiedene Konstante $c - c'$ durch $\pi$ teilbar wäre.

6. \gesperrt{Jede ganze Funktion $\omega$ in $\Omega$ ist nach dem Modul $\pi$ mit einer Konstante $b$ kongruent, d.\,h. man kann die Konstante $b$ so wählen, das
$\omega - b$ durch $\pi$ teilbar wird.}
% [sic] source prints "das" for "daß".

Denn die Funktion $\omega$ genügt nach \S\,175 einer Gleichung:
\[
\omega^n + a_1 \omega^{n-1} + \cdots + a_{n-1}\omega + a_n = 0,
\]
worin die $a_1, \ldots, a_n$ ganze rationale Funktionen von $z$ sind. Sind $a_1^0, \ldots, a_n^0$ die Reste dieser Funktionen bei der Teilung durch $z - c$, also Konstanten, so folgt:
\[
\omega^n + a_1^0 \omega^{n-1} + \cdots + a_{n-1}^0 \omega + a_n^0 \equiv 0\ (\mathrm{mod}\ \pi),
\]
und wenn man die Funktion auf der linken Seite in Linearfaktoren zerlegt:
\weq{1}{(\omega - b)(\omega - b')(\omega - b'') \cdots \equiv 0\ (\mathrm{mod}\ \pi).}
Es muß also wenigstens einer dieser Linearfaktoren durch $\pi$ teilbar sein, was zu beweisen war.

Aus der nachgewiesenen einwertigen Zerlegbarkeit der Funktionale in Primfaktoren ergeben sich weitgehende Folgerungen, von denen die wichtigsten hier angeführt werden sollen.

Wenn in einem Funktional $\varphi$ des Körpers $\Omega$ eine gewisse Anzahl der Hilfsvariablen, $x, y, \ldots$, nur im Zähler vorkommen, so möge
\[
\varphi = \varphi(x, y, \ldots)
\]
eine \gesperrt{holomorphe Funktion} von $x, y, \ldots$ heißen. Von denen gilt der Satz:

7. \gesperrt{Eine holomorphe Funktion ist nur dann ein ganzes Funktional in $\overline{\Omega}$, wenn die Koeffizienten der geordneten Funktion $\varphi$ ganze Funktionale sind.}

Es genügt offenbar, den Satz für holomorphe Funktionen einer Variablen
\weq{2}{\varphi = x^m \varphi_0 + x^{m-1}\varphi_1 + \cdots + x \varphi_{m-1} + \varphi_m}
nachzuweisen, weil er daraus durch vollständige Induktion allgemein bewiesen werden kann.

Dieser Beweis ergibt sich aber aus dem \gesperrt{Gaussschen} Satz:
%CONT>>
